% This work is made available under the terms of the
% Creative Commons Attribution-ShareAlike 4.0 license,
% http://creativecommons.org/licenses/by-sa/4.0/.
%
% Version: $Revision: 4890 $

\documentclass[a4paper]{book}

\usepackage{wrapfig}
\usepackage{graphicx}
\usepackage{hyperref}
\usepackage{multirow}
\usepackage{scalefnt}
\usepackage{tikz}

% watermark -- for draft stage
\usepackage[firstpage]{draftwatermark}
\SetWatermarkLightness{0.9}
\SetWatermarkScale{5}

% Copyright (c) 2009 by the University of Waikato, Hamilton, NZ. 
% This work is made available under the terms of the 
% Creative Commons Attribution-ShareAlike 3.0 license, 
% http://creativecommons.org/licenses/by-sa/3.0/. 
%
% Version: $Revision$

\newenvironment{tight_itemize}{
\begin{itemize}
  \setlength{\itemsep}{1pt}
  \setlength{\parskip}{0pt}
  \setlength{\parsep}{0pt}}{\end{itemize}
}

\newenvironment{tight_enumerate}{
\begin{enumerate}
  \setlength{\itemsep}{1pt}
  \setlength{\parskip}{0pt}
  \setlength{\parsep}{0pt}}{\end{enumerate}
}

% if you just need a simple heading
% Usage:
%   \heading{the text of the heading}
\newcommand{\heading}[1]{
  \vspace{0.3cm} \noindent \textbf{#1} \newline
}

\newcommand{\icon}[1]{\tikz[baseline=-3pt]\node[inner sep=0pt,outer sep=0pt]{\includegraphics[height=1.1em]{#1}};}


\title{
  \textbf{ADAMS} \\
  {\Large \textbf{A}dvanced \textbf{D}ata mining \textbf{A}nd \textbf{M}achine
  learning \textbf{S}ystem} \\
  {\Large Module: adams-compress} \\
  \vspace{1cm}
  \includegraphics[width=2cm]{images/compress-module.png} \\
}
\author{
  Peter Reutemann
}

\setcounter{secnumdepth}{3}
\setcounter{tocdepth}{3}

\begin{document}

\begin{titlepage}
\maketitle

\thispagestyle{empty}
\center
\begin{table}[b]
	\begin{tabular}{c l l}
		\parbox[c][2cm]{2cm}{\copyright 2012-2017} &
		\parbox[c][2cm]{5cm}{\includegraphics[width=5cm]{images/coat_of_arms.pdf}} \\
	\end{tabular}
	\includegraphics[width=12cm]{images/cc.png} \\
\end{table}

\end{titlepage}

\tableofcontents
%\listoffigures
%\listoftables

%%%%%%%%%%%%%%%%%%%%%%%%%%%%%%%%%%%
\chapter{Introduction}
Compressing and decompressing files, archiving multiple files, extracting files
from archives. These are all actions that are very common and most of the time
manually performed. Using ADAMS, these steps can be automated thanks to the 
various transformers available that handle archives.

In general, there are two different kinds of actors:
\begin{tight_itemize}
  \item multi-file archives (compress/decompress)
  \item single-file algorithms (compression/decompression)
\end{tight_itemize}
The latter only work on a single file, like for instance gzip. ZIP, on the
other hand, works with multiple files.

The following chapters cover the actors for compression and decompression in 
more detail.

%%%%%%%%%%%%%%%%%%%%%%%%%%%%%%%%%%%
\chapter{Compression}
In order to compress one or more files, you can use the following transformers:
\begin{tight_itemize}
	\item \textbf{single files}
	\begin{tight_itemize}
		\item Bzip2 (see \cite{bzip2})
		\item GZIP (see \cite{gzip})
		\item Lzf (see \cite{lzf})
		\item Lzma (see \cite{lzma})
	\end{tight_itemize}
	\item \textbf{multiple files}
	\begin{tight_itemize}
		\item Tar (see \cite{tar})
		\item ZIP (see \cite{zip})
	\end{tight_itemize}
\end{tight_itemize}

\heading{Single file}
These transformers take a single file as input, which gets 
compressed.\footnote{adams-compress-gzip\_single\_file.flow}
You have the choice of selecting a target file. If not, then the generated
archive gets placed in the same directory as the input file. Optionally,
you can also remove the original file. This is useful if you simply want
to save disk-space and compress all files in a directory, but don't need
the original files anymore.

\heading{Multiple files}
Transformers that manage archives instead of only compressing single files,
you need to supply them with an array of file names that should get added
to the archive (no incremental adding possible at the moment). Using the
\textit{stripPath} regular expression, you can exert control over what of
the (most likely absolute) path of the file names should end up in the
archive. In order to strip the complete path from the names, simply use
``.*'' as expression.\footnote{adams-compress-generate\_zip.flow}


%%%%%%%%%%%%%%%%%%%%%%%%%%%%%%%%%%%
\chapter{Decompression}
In order to decompress a compressed file or extract one more files from a
multi-file archive, you have the following transformers available:
\begin{tight_itemize}
	\item \textbf{single file archives}
	\begin{tight_itemize}
		\item UnBzip2 (see \cite{bzip2})
		\item UnGZIP (see \cite{gzip})
		\item UnLzf (see \cite{lzf})
		\item UnLzma (see \cite{lzma})
	\end{tight_itemize}
	\item \textbf{multi-file archives}
	\begin{tight_itemize}
		\item UnRAR (see \cite{rar})
		\item UnTar (see \cite{tar})
		\item UnZIP (see \cite{zip})
	\end{tight_itemize}
\end{tight_itemize}

\heading{Single file archives}
When using transformers for decompressing the content of a single file archive,
you can choose whether you would like to extract the content to a different 
directory and even whether to use a different file name (by default, the file name
is the one with the compression's suffix).\footnote{adams-compress-decompress\_gzipped\_file.flow}

\heading{Multi-file archives}
By default, all files from a multi-file archive get extracted. If you only 
require a subset of the files, you can use the \textit{regExp} option to
limit the files extracted to the ones that match the regular expression (you
can also invert the matching sense).\footnote{adams-compress-extract\_from\_zip.flow}
Optionally, you can choose whether to re-create the directory structure 
stored in the archive.

%%%%%%%%%%%%%%%%%%%%%%%%%%%%%%%%%%%
% Copyright (c) 2009-2012 by the University of Waikato, Hamilton, NZ. 
% This work is made available under the terms of the 
% Creative Commons Attribution-ShareAlike 4.0 license,
% http://creativecommons.org/licenses/by-sa/4.0/.
%
% Version: $Revision$

\begin{thebibliography}{999}
	% to make the bibliography appear in the TOC
	\addcontentsline{toc}{chapter}{Bibliography}

    % references
	\bibitem{adams}
		\textit{ADAMS} -- Advanced Data mining and Machine learning System \\
		\url{https://adams.cms.waikato.ac.nz/}{}

	\bibitem{esrigrid}
	 	\textit{Esri Grid} -- a raster GIS file format deveoped by Esri. \\
		\url{https://en.wikipedia.org/wiki/Esri\_grid}{}

	\bibitem{kml}
	 	\textit{Keyhole Markup Language} -- an XML notation for expressing
	 	geographic annotation and visualization within Internet-based,
	 	two-dimensional maps and three-dimensional Earth browsers. \\
		\url{http://en.wikipedia.org/wiki/Keyhole\_Markup\_Language}{}

	\bibitem{postgresql}
	 	\textit{PostgreSQL} -- a powerful, open source object-relational
	 	database system. \\
		\url{http://www.postgresql.org/}{}

	\bibitem{postgis}
		\textit{PostGIS} -- a spatial database extender for PostgreSQL
		object-relational database. It adds support for geographic
		objects allowing location queries to be run in SQL.  \\
		\url{http://postgis.net/}{}

	\bibitem{srid4269}
	 	\textit{SRID 4269} -- or NAD 83 (North American Datum). \\
		\url{http://spatialreference.org/ref/epsg/4269/}{}

	\bibitem{mysql}
		\textit{MySQL} -- an open-source relational database management
		system (RDBMS) \\
		\url{http://www.mysql.com/}{}

\end{thebibliography}


\end{document}
